\newpage
\thispagestyle{empty}

\begin{par}
\setlength{\baselineskip}{5mm}

\begin{center}
    \textbf{АННОТАЦИЯ} \\
    \vspace{5mm}
    ВЫПУСКНАЯ КВАЛИФИКАЦИОННАЯ РАБОТА БАКАЛАВРА \\
    \noindent
    \begin{flushleft}
        Наименование темы \uline{\parbox[t]{\linewidth}{\topic\hfill}} \\
    \hrulefill
    \end{flushleft}
\end{center}

\begin{flushleft}
    Выполнена студентом(кой) \uline{\studentforabstract\hfill} \\
    \faculty, \university \\
    Кафедра \uline{\department\hfill} \\
    Группа \uline{\group} \\
    Направление подготовки: \programnum\thinspace\programname \\
    Направленность (профиль): \specialization \\
    \vspace{5mm}
    Объем работы: \total{page} страниц \\
    Количество иллюстраций: \totalfigures \\
    Количество таблиц: \totaltables \\
    Количество литературных источников: \total{citenum} \\
    Количество приложений: \total{appendixnum} \\
    Ключевые слова: биосонификация, генерация музыки, нейронные сети, MIDI, FASTA, трансформер, символическая музыка, биологические последовательности, POP909, глубокое обучение
\end{flushleft}

В работе рассматривается задача музыкальной сонификации биологических данных. Объектом исследования являются биологические данные и методы их преобразования в музыкальную форму. Цель работы — разработать и экспериментально проверить самообучающуюся модель искусственного интеллекта для синтеза и оценки мелодий по биологическим данным.

Для достижения цели изучены подходы к биосонификации и генерации символической музыки, подготовлены биологические и музыкальные данные, реализовано извлечение признаков из DNA, RNA и protein последовательностей, разработана иерархическая архитектура генерации гармонии и мелодии, обучены модели на доступном оборудовании и проведена экспериментальная оценка.

Разработана система BioSonification, реализующая иерархическую генерацию музыки Bio→Harmony→Melody. Система обучена на корпусе POP909 и геномных данных RefSeq. Актуальная модель веб-интерфейса показала validation loss 0.1454 для гармонии и 0.1566 для мелодии. Экспериментальная оценка подтвердила, что сгенерированные композиции являются валидными двухдорожечными MIDI-файлами, а мелодии согласуются с гармонией: chord-tone ratio составил 0.6710 против 0.4043 у случайного baseline и 0.6073 у предыдущей версии модели.

Научная новизна состоит в использовании биологических признаков как управляющего сигнала для структурированной музыкальной генерации, а не как прямого отображения символов в ноты. Практическая значимость работы заключается в создании рабочего программного прототипа с веб-интерфейсом, обученными моделями и воспроизводимыми результатами. Система может применяться в биоинформатике, научной коммуникации и образовании.

\begin{flushleft}
    \vspace{1.02cm}
    \vfill
    \student, \\
    \readydate\space\currentyear
\end{flushleft}

\end{par}
